% Existing questions and full draft answers retained for reference.
\section*{Research questions and answers}
\addcontentsline{toc}{section}{Research questions and answers}
\label{sec:working-research-questions}

\subsection*{RQ1 --- Can fine-grained topics be discovered with nothing labelled?}
Can an unsupervised pipeline recover subject structure at the scale of a real web corpus,
and what does it get wrong? (Section~\ref{sec:rq1})

\paragraph{Answer.}
\gd{Yes, with two qualifications.} Fine-grained subject topics can be recovered from an
unlabelled bilingual web corpus at the scale of 616{,}499 documents, and every document can be
assigned to one. But the corpus has no natural number of topics, so the resolution must be
chosen on grounds other than geometry (RQ3); and roughly a fifth to a quarter of the
discovered topics describe page roles rather than subjects, a failure the standard coherence
metrics cannot detect.

\subsection*{RQ2 --- Which feature should be embedded?}
A web page offers body text, a title, a description, and a URL. Which of these, alone or in
combination, should be given to the embedding model? (Section~\ref{sec:rq2})

\paragraph{Answer.}
\gd{Plain body text}, not because it produces the most coherent topics, which it does not in
German, but because it is the only input present on every document and the only one that does
not inherit a particular site's template conventions. \measured{} No variant reliably improves
topic quality; \proved{} every alternative variant forfeits guaranteed coverage.

\subsection*{RQ3 --- How much granularity is enough?}
Topic models must be told, directly or indirectly, how fine to go. Is there a principled way
to decide, and does that way transfer to corpora it was not fitted on? (Sections~\ref{sec:rq3}
and~\ref{sec:neutrality})

\paragraph{Answer.}
\gd{Enough granularity is the setting that maximises the number of coarse topics that are real
splits}, where a coarse topic is a real split when it has at least two fine children \emph{and} no
two of those children are near-duplicates of each other. The first condition rules out renames;
the second rules out divisions that are nominal. \measured{} On the motivating corpus this selects
\textbf{min 100 in both languages}, 9 of 20 in English and 10 of 20 in German.

The criterion needs no coherence metric and no human judgement. The threshold of two is
definitional. The duplicate cutoff is a free parameter but is shown not to be load-bearing: the
selected setting is identical for every cutoff from $0.34$ to $1.00$, and at $1.00$ the rule reduces to ``no two children have identical
keyword lists'', which introduces no constant at all. Because the count is not monotone in the
sweep direction, it identifies a setting rather than merely preferring the finest one available.

It detects two failure modes that coherence-based selection cannot see: whole themes silently
losing their fine representation, and fine topics that subdivide a theme in name only.

